% The font could be set to Windows-specific Calibri by using the 'calibri' option
% The default is TeX Gyre Heros (metric-compatible with Arial)
\documentclass[]{mcdowellcv}

% For mathematical symbols
\usepackage{amsmath}

% For hyperlinking
\usepackage[hidelinks]{hyperref}

% Set applicant's personal data for header
\name{Corben Dredge}
\contacts{\href{tel:07*********}{07*********} \linebreak \href{mailto:cdredge@protonmail.ch}{cdredge@protonmail.ch} \linebreak \href{https://corbenshiro.github.io/}{corbenshiro.github.io/}}

% Where the magic happens
\begin{document}

	% Print the header
	\makeheader

	% Print the content
	\begin{cvsection}{Employment}
		\begin{cvsubsection}{Software Engineer, Intern}{Red Hat, Inc.}{May - August 2019}
			\begin{itemize}
				\item Worked on RPM packaging through spec file engineering.
				\item Bootstrapped (via Podman) containerized Apache Maven and Gradle environments with the Supersonic Subatomic Java implementation known as Quarkus.
				\item Tinkered with Unix-like systems such as GNU/Linux (various distributions), Alpine Linux (OpenRC with BusyBox compiled against musl libc) and FreeBSD whilst conversing about them with fellow colleagues.
			\end{itemize}
		\end{cvsubsection}
		\begin{cvsubsection}{IT Staff Assistant, Intern}{Newham Sixth Form College}{December 2017}
			\begin{itemize}
				\item Rebuilt workstations and flashed its storage drives with operating system disk images (mainly Windows 7 Professional and Windows 10 Education editions) for later use within the college.
				\item Fixed printing errors via a reload of all the printers' respective daemons.
			\end{itemize}
		\end{cvsubsection}
	\end{cvsection}

	\begin{cvsection}{Education}
		\begin{cvsubsection}{Computer Science}{Newham Sixth Form College}{September 2018 - July 2019}
			\begin{itemize}
				\item Pearson BTEC Level 3 90-credit Diploma in Information Technology
			\end{itemize}
		\end{cvsubsection}

		\begin{cvsubsection}{Networking and Systems}{Newham Sixth Form College}{September 2017 - July 2018}
			\begin{itemize}
				\item Pearson BTEC Level 2 First Extended Certificate in Information and Creative Technology
			\end{itemize}
		\end{cvsubsection}
	\end{cvsection}

	\begin{cvsection}{Personal Summary}
		\begin{cvsubsection}{}{}{}
			\begin{itemize}
				\item I'm familiar with both ends of conversation — speaking and listening, I have become accustomed to developing rapport as well as being affable.
				\item I'm very thorough in regard to my interests in IT, therefore I've acquired many skills that revolve around computing (including subcategories of cybersecurity, software engineering and technical writing).
			\end{itemize}
		\end{cvsubsection}
		\begin{cvsubsection}{Verified Boot}{}{}	
			\begin{itemize}
				\item Cryptographically utilises hardware-backed root of trust in order to eliminate malware persistence.
				\begin{itemize}
					\item Android (AOSP)
					\begin{itemize}
						\item GrapheneOS
					\end{itemize}
					\item Chrome OS (Chromium OS)
					\item iOS/iPadOS
					\item macOS
					\item Windows NT (preferably Windows 10 on a Secured-core PC)
					\item Qubes OS (with the Anti Evil Maid module explicitly installed and enabled)
				\end{itemize}
			\end{itemize}
		\end{cvsubsection}
		\begin{cvsubsection}{Mandatory Access Control}{}{}
			\begin{itemize}
				\item Infinitely preferable to DAC; full system MAC policies are highly encouraged.
				\begin{itemize}
					\item SELinux
					\begin{itemize}
						\item AOSP and by inheritance; GrapheneOS's stricter base SELinux policy.
						\item Chromium OS
					\end{itemize}
					\item AppArmor
					\begin{itemize}
						\item apparmor-profile-everything
						\item sandbox-app-launcher (requires bubblewrap)
					\end{itemize}
					\item System Integrity Protection (macOS)
					\item Mandatory Integrity Control (Windows NT)
				\end{itemize}
			\end{itemize}
		\end{cvsubsection}
		\begin{cvsubsection}{Full-disk Encryption}{}{}
			\begin{itemize}
				\item Recommended to use in conjunction with verified boot as it cannot safeguard boot by itself.
				\begin{itemize}
					\item LUKS (dm-crypt)
					\item VeraCrypt
					\item FileVault 2
					\item BitLocker
				\end{itemize}
			\end{itemize}
		\end{cvsubsection}
		\begin{cvsubsection}{Memory-safe Langages}{}{}
			\begin{itemize}
				\item A subset of recommended practices in memory safety that holistically pertains to programming.
				\begin{itemize}
					\item Java/Kotlin (Android favours the latter for application development)
					\item Swift (developed by Apple for their operating systems; intended as a memory-safe (among other things) replacement for Objective-C)
					\item Rust (Windows NT is adopting this)
					\item Go (the all-in-one — statically typed, compiled programming language with memory safety, garbage collection, structural typing and CSP-style concurrency)
				\end{itemize}
			\end{itemize}
		\end{cvsubsection}
		\begin{cvsubsection}{Virtualization-based Security}{}{}
			\begin{itemize}
				\item The Windows 10 host becomes its own VM, thus ensuring memory integrity via Hypervisor-Enforced Code Integrity.
				\begin{itemize}
					\item TPM 2.0, UEFI and Hyper-V are prerequisites.
				\end{itemize}
			\end{itemize}
		\end{cvsubsection}
		\begin{cvsubsection}{Xen-based Virtualization}{}{}
			\begin{itemize}
				\item Security by compartimentalization.
				\begin{itemize}
					\item Qubes OS
					\begin{itemize}
						\item Split GPG (preferable to utilise vault for Master Keyring as well as to allocate a network-isolated VM each with their own respective subkey)
						\item Split SSH (preferable to utilise vault via KeePassXC with SSH Agent integration)
					\end{itemize}
				\end{itemize}
			\end{itemize}
		\end{cvsubsection}
		\begin{cvsubsection}{STIGs}{}{}
			\begin{itemize}
				\item Windows 10 V2R1
				\item Windows Server 2019 V2R1
				\item Windows Defender Antivirus V2R1
				\item Windows Firewall V1R7
				\item Microsoft Office 2019/Office 365 Pro Plus V1R2
				\item Oracle JRE 8 V1R5
			\end{itemize}
		\end{cvsubsection}
		\begin{cvsubsection}{Development Environments}{}{}
			\begin{itemize}
				\item MSYS2 package groups
				\begin{itemize}
					\item base-devel
					\item mingw-w64-x86\_64-toolchain
				\end{itemize}
				\item Revision control systems
				\begin{itemize}
					\item Git (command-line reference is stylised as git)
				\end{itemize}
			\end{itemize}
		\end{cvsubsection}
		\begin{cvsubsection}{Word Processors}{}{}
			\begin{itemize}
				\item Microsoft Word
				\item LibreOffice Writer
			\end{itemize}
		\end{cvsubsection}
		\begin{cvsubsection}{Typesetting Automation}{}{}
			\begin{itemize}
				\item \TeX
				\begin{itemize}
					\item \LaTeX
					\begin{itemize}
						\item Lua\LaTeX
					\end{itemize}
				\end{itemize}
			\end{itemize}
		\end{cvsubsection}
	\end{cvsection}

	\begin{cvsection}{Authenticity}
		\begin{cvsubsection}{}{}{}
			\begin{itemize}
				\item For details unto how I secure my website, webpages and documents (including this curriculum vitæ); refer to:
				\begin{itemize}
					\item \href{https://corbenshiro.github.io/security-of-this-site/}{corbenshiro.github.io/security-of-this-site/}
				\end{itemize}
			\end{itemize}
		\end{cvsubsection}
	\end{cvsection}

	\begin{cvsection}{Extracurricular Interests}
		\begin{cvsubsection}{Sports and Fitness}{}{}
			\begin{itemize}
				\item Badminton
				\item Weightlifting
			\end{itemize}
		\end{cvsubsection}
		\begin{cvsubsection}{Franchises}{}{}
			\begin{itemize}
				\item Vanpaia Hantā Dī
				\item Hokuto no Ken
				\item Warhammer 40K
				\item Freedom Force
			\end{itemize}
		\end{cvsubsection}
		\begin{cvsubsection}{Geek Code (v3.12)}{}{}
			\begin{itemize}
				\item GCS/IT/TW !d s+:- a? C U++ P+ K- w++++ Y+ PGP+ R* !tv h++ !y+
			\end{itemize}
		\end{cvsubsection}
	\end{cvsection}

	\begin{cvsection}{References}
		\begin{cvsubsection}{}{}{}
			\begin{itemize}
				\item They are available upon request.
			\end{itemize}
		\end{cvsubsection}
	\end{cvsection}

\end{document}
